% ===========================================================================
%  Capítulo 4 — Generación de código
% ===========================================================================
\chapter{Generaci\'on de c\'odigo}
\label{cap:generacion}

La generaci\'on de c\'odigo es la etapa del sistema que traduce el AST
construido por las formas del DSL a un programa Python que produce el archivo
\xlsx{}.  Este cap\'{\i}tulo describe c\'omo el backend resuelve las
coordenadas Excel, c\'omo compila cada tipo de nodo de expresi\'on a una
cadena de f\'ormula, y c\'omo genera las f\'ormulas SUMPRODUCT que
implementan el formato condicional persistente.

% ---------------------------------------------------------------------------
\section{Arquitectura del backend}
\label{sec:arquitectura-backend}
% ---------------------------------------------------------------------------

El backend est\'a organizado en dos archivos con responsabilidades separadas:

\begin{description}
  \item[\file{code-generation-utils.lisp}.] Funciones auxiliares, variables
    din\'amicas de configuraci\'on y la funci\'on de entrada principal.
    Incluye: \dsl{col->letter} (conversi\'on de n\'umero de columna a letra
    Excel), \dsl{build-col-map} (construcci\'on del mapa s\'{\i}mbolo~DSL
    $\to$ letra), \dsl{xl-generate} y \dsl{xl-run-generated}.

  \item[\file{generate-code-direct.lisp}.] Los m\'etodos CLOS de la
    funci\'on gen\'erica \dsl{compile-excel-formula} y los m\'etodos del
    generador de c\'odigo Python.  Este archivo es el \'unico lugar del
    sistema donde pueden aparecer letras de columna Excel, cadenas de
    f\'ormula y llamadas a openpyxl.
\end{description}

El principio de separaci\'on de responsabilidades queda codificado en la
propia arquitectura: \file{ast-def.lisp} declara \emph{qu\'e existe},
\file{dsl-directo.lisp} declara \emph{c\'omo se escribe} y
\file{generate-code-direct.lisp} declara \emph{a qu\'e se traduce}.

% ---------------------------------------------------------------------------
\section{Resoluci\'on de coordenadas}
\label{sec:resolucion-coordenadas}
% ---------------------------------------------------------------------------

% ............................................................................
\subsection{Conversi\'on de n\'umero de columna a letra Excel}
\label{ssec:col-letter}
% ............................................................................

La funci\'on \dsl{col->letter} convierte un entero $n$ (la posici\'on
ordinal de una columna, con base~1) a la letra o letras de columna
correspondientes en Excel: columna~1~$\to$~\texttt{A}, columna~26~$\to$~\texttt{Z},
columna~27~$\to$~\texttt{AA}, etc.

\begin{lstlisting}[language=CommonLisp,
  caption={Funci\'on col-\textgreater{}letter.}]
(defun col->letter (n)
  (with-output-to-string (s)
    (loop while (> n 0)
          do (multiple-value-bind (q r) (floor (1- n) 26)
               (setf n q)
               (write-char (code-char (+ 65 r)) s)))))
\end{lstlisting}

El algoritmo es una conversi\'on de base 10 a base 26 con el ajuste
habitual ($\lfloor (n-1)/26 \rfloor$, $((n-1) \bmod 26) + 65$) que mapea
los residuos a los caracteres ASCII de las letras may\'usculas.

% ............................................................................
\subsection{Construcci\'on del mapa de columnas}
\label{ssec:col-map}
% ............................................................................

Para cada tabla de la regi\'on se construye un \emph{mapa de columnas}:
una lista de asociaci\'on que relaciona cada s\'{\i}mbolo DSL de columna con
su letra Excel absoluta.  La funci\'on \dsl{build-col-map} construye el mapa
de una tabla sola, asumiendo que empieza en la columna~1.  La funci\'on
\dsl{build-table-col-maps} construye el mapa global de toda la regi\'on,
teniendo en cuenta el desplazamiento de cada tabla dentro de la regi\'on:

\begin{lstlisting}[language=CommonLisp,
  caption={Construcci\'on del mapa de columnas de una regi\'on.}]
(defun build-col-map (tbl)
  (loop for name in (col-names tbl)
        for i from 1
        collect (cons name (col->letter i))))

(defun build-table-col-maps (tables offsets)
  (loop for tbl in tables
        for offset in offsets
        collect (cons (id tbl)
                      (loop for name in (col-names tbl)
                            for i from 1
                            collect (cons name
                                          (col->letter (+ offset i)))))))
\end{lstlisting}

El resultado se almacena en la variable din\'amica \dsl{*table-col-maps*},
disponible durante toda la compilaci\'on de la regi\'on.

% ............................................................................
\subsection{Cálculo de rangos de fila}
\label{ssec:rangos-fila}
% ............................................................................

Cada tabla tiene una primera fila de datos (\dsl{first-row}, configurable,
por defecto~4) y una \'ultima fila que depende del n\'umero de filas de datos
y de \dsl{cell-height}:

\begin{equation}
\textit{last-row} = \textit{first-row} + |\textit{datos}| \times \textit{cell-height} - 1
\end{equation}

Estos valores se almacenan en la variable din\'amica \dsl{*table-data-rows*}
como una lista de asociaci\'on \dsl{(table-id . (first-row . last-row))},
lo que permite que los nodos \dsl{xl-table-range} usen los rangos de fila
correctos para cada tabla aunque la compilaci\'on se realice dentro del
contexto de una tabla diferente.

% ---------------------------------------------------------------------------
\section{La funci\'on gen\'erica \texttt{compile-excel-formula}}
\label{sec:compile-excel-formula}
% ---------------------------------------------------------------------------

El m\'etodo central del backend es la funci\'on gen\'erica
\dsl{compile-excel-formula}, que toma un nodo AST y devuelve la cadena de
f\'ormula Excel correspondiente.  Su firma es:

\begin{lstlisting}[language=CommonLisp,
  caption={Declaraci\'on de compile-excel-formula.}]
(defgeneric compile-excel-formula
    (expr col-map data-names row-num first-row last-row)
  (:documentation
   "Compila un nodo AST de expresion a string de formula Excel"))
\end{lstlisting}

Los par\'ametros son:
\begin{description}
  \item[\dsl{expr}] El nodo AST a compilar.
  \item[\dsl{col-map}] El mapa de columnas de la tabla actual
    (\dsl{sym $\to$ letra}).
  \item[\dsl{data-names}] Lista de s\'{\i}mbolos DSL de la tabla actual
    (para \dsl{VLOOKUP}).
  \item[\dsl{row-num}] N\'umero de fila Excel que se est\'a compilando
    (para referencias relativas).
  \item[\dsl{first-row}] Primera fila de datos de la tabla actual.
  \item[\dsl{last-row}] \'Ultima fila de datos de la tabla actual.
\end{description}

El despacho m\'ultiple de CLOS garantiza que, para cada tipo de nodo, se
ejecuta el m\'etodo correcto sin ning\'un \dsl{cond} centralizado.
A\~nadir un nuevo tipo de expresi\'on al lenguaje s\'olo requiere:
\begin{enumerate}
  \item Definir una nueva clase con \dsl{defclass*} en \file{ast-def.lisp}.
  \item Definir un macro en \file{dsl-directo.lisp} que construya el nodo.
  \item Definir un \dsl{defmethod} de \dsl{compile-excel-formula} en
    \file{generate-code-direct.lisp}.
\end{enumerate}

Ning\'un c\'odigo existente necesita modificarse.

% ............................................................................
\subsection{M\'etodos para expresiones simples}
\label{ssec:metodos-simples}
% ............................................................................

\begin{lstlisting}[language=CommonLisp,
  caption={M\'etodos compile-excel-formula para expresiones simples.}]
;; Referencia a columna de la tabla actual: "K4"
(defmethod compile-excel-formula
    ((e clase-xl-expr-column-ref) col-map data-names
     row-num first-row last-row)
  (let* ((letter (cdr (assoc (name e) col-map)))
         (ctx (context e)))
    (if (typep ctx 'clase-xl-expr-previous-row)
        (format nil "~a~a" letter (1- row-num))
        (format nil "~a~a" letter row-num))))

;; Comparacion de igualdad: "K4=B4"
(defmethod compile-excel-formula
    ((e clase-xl-expr-equals) col-map data-names
     row-num first-row last-row)
  (format nil "~a=~a"
          (compile-excel-formula (a e) col-map data-names
                                  row-num first-row last-row)
          (compile-excel-formula (b e) col-map data-names
                                  row-num first-row last-row)))

;; Literal de texto: "\"Total\""
(defmethod compile-excel-formula
    ((e clase-xl-expr-string) col-map data-names
     row-num first-row last-row)
  (format nil "~s" (value e)))
\end{lstlisting}

% ............................................................................
\subsection{M\'etodos para rangos y agregados}
\label{ssec:metodos-rangos}
% ............................................................................

\begin{lstlisting}[language=CommonLisp,
  caption={M\'etodo para xl-table-range.}]
;; Rango absoluto de columnas de una tabla especifica de la region.
;; Usa *table-data-rows* para obtener los limites exactos de esa tabla.
(defmethod compile-excel-formula
    ((e clase-xl-table-range) col-map data-names
     row-num first-row last-row)
  (let* ((tmap        (cdr (assoc (table-id e) *table-col-maps*)))
         (from-letter (cdr (assoc (name (from-col e)) tmap)))
         (to-letter   (if (to-col e)
                          (cdr (assoc (name (to-col e)) tmap))
                          from-letter))
         (tbl-rows    (cdr (assoc (table-id e) *table-data-rows*)))
         (act-first   (if tbl-rows (car tbl-rows) first-row))
         (act-last    (if tbl-rows (cdr tbl-rows) last-row)))
    (format nil "$~a$~a:$~a$~a"
            from-letter act-first to-letter act-last)))
\end{lstlisting}

La clave t\'ecnica de este m\'etodo es que consulta \dsl{*table-data-rows*}
con el identificador de la \emph{tabla referenciada}, no con el de la tabla
actual.  As\'{\i}, \dsl{(trange defensas-table tutor secretario)} compila al
rango de \dsl{defensas-table} aunque se est\'e compilando una f\'ormula de
\dsl{profesores-table}.

% ---------------------------------------------------------------------------
\section{Generaci\'on de formato condicional}
\label{sec:cf-generacion}
% ---------------------------------------------------------------------------

El formato condicional es el aspecto t\'ecnicamente m\'as elaborado de la
generaci\'on de c\'odigo.  El backend recibe un nodo \dsl{xl-expr-exists} y
produce una f\'ormula SUMPRODUCT que implementa la sem\'antica del
cuantificador existencial declarado.

% ............................................................................
\subsection{Variante same-table: SUMPRODUCT con exclusi\'on de fila actual}
\label{ssec:cf-same-table}
% ............................................................................

Cuando \dsl{exists} usa \dsl{:all-rows}, el backend genera una f\'ormula
SUMPRODUCT que, para cada fila, multiplica tres factores:
\begin{enumerate}
  \item Coincidencia en las columnas de agrupaci\'on (\dsl{:matching}).
  \item Exclusi\'on de la fila actual (para no comparar una fila consigo misma).
  \item Solapamiento en al menos una columna de rol (\dsl{:any-overlap}).
\end{enumerate}

Para una regla con \dsl{:matching (dia hora)} y
\dsl{:any-overlap (tutor \ldots\ secretario)}, la f\'ormula generada para
la fila~4 es:

\begin{lstlisting}[style=pythonstyle,
  caption={F\'ormula Excel CF para la variante same-table (fila 4).}]
=SUMPRODUCT(
    ($G$4:$G$6=$G4) *          -- dia coincide en ambas filas
    ($H$4:$H$6=$H4) *          -- hora coincide en ambas filas
    (ROW($B$4:$B$6)<>ROW($B4)) *  -- excluir la fila actual
    (($B$4:$B$6=$B4)+($C$4:$C$6=$C4)+...+($F$4:$F$6=$F4))
) > 0
\end{lstlisting}

Las referencias con columna fija y fila relativa (p.\,ej.\ \texttt{\$G4})
hacen que la f\'ormula se eval\'ue correctamente para cada fila del rango:
al desplazarse de la fila~4 a la fila~5, las referencias de ``fila actual''
avanzan autom\'aticamente.  El resultado de \texttt{SUMPRODUCT} es el n\'umero
de filas del dominio que satisfacen las tres condiciones; la comparaci\'on
\texttt{> 0} convierte ese conteo en un valor booleano para la regla de
formato condicional.

% ............................................................................
\subsection{Variante cross-table: SUMPRODUCT con COUNTIFS}
\label{ssec:cf-cross-table}
% ............................................................................

Cuando \dsl{exists} usa \dsl{:rows-of}, el backend genera una f\'ormula
SUMPRODUCT que combina dos condiciones:
\begin{enumerate}
  \item Que el valor de la celda actual aparezca en al menos una columna de
    rol de la tabla referenciada (\dsl{:self-in}).
  \item Que el slot de tiempo de esa fila tenga m\'as de una defensa
    programada (\dsl{:matching} con \texttt{COUNTIFS}).
\end{enumerate}

Para la regla de \dsl{profesores-table} con
\dsl{:self-in (tutor \ldots\ secretario)} y \dsl{:matching (dia hora)},
la f\'ormula generada para la celda K4 es:

\begin{lstlisting}[style=pythonstyle,
  caption={F\'ormula Excel CF para la variante cross-table (celda K4).}]
=SUMPRODUCT(
    (($B$4:$B$6=$K4)+($C$4:$C$6=$K4)+($D$4:$D$6=$K4)+
     ($E$4:$E$6=$K4)+($F$4:$F$6=$K4)) *
    (COUNTIFS($G$4:$G$6,$G$4:$G$6,
              $H$4:$H$6,$H$4:$H$6) > 1)
) > 0
\end{lstlisting}

La referencia \texttt{\$K4} corresponde al nombre del profesor en la celda
actual (columna~K bloqueada, fila relativa).  Los rangos
\texttt{\$B\$4:\$B\$6} a \texttt{\$F\$4:\$F\$6} son las columnas de rol de
\dsl{defensas-table}, completamente absolutas.

\texttt{COUNTIFS} con cada par de columnas de coincidencia iguales a s\'{\i}
mismas produce el n\'umero de filas que comparten exactamente los mismos
valores en esas columnas; comparar el resultado con~$1$ selecciona los
slots de tiempo con m\'as de una defensa.  El producto SUMPRODUCT de ambas
condiciones devuelve el n\'umero de filas de \dsl{defensas-table} en que el
profesor actual participa en un slot conflictivo.

Al ser una \texttt{FormulaRule} incrustada en el archivo Excel, cualquier
edici\'on posterior de los datos provoca una reevaluaci\'on autom\'atica del
coloreado, sin necesidad de volver a ejecutar el DSL.

% ---------------------------------------------------------------------------
\section{Generaci\'on de referencias cruzadas entre hojas}
\label{sec:cf-cross-sheet}
% ---------------------------------------------------------------------------

El nodo \dsl{xl-expr-collect-over} itera sobre una lista de hojas del
workbook y agrega la expresi\'on del cuerpo para cada una.  El mecanismo de
iteraci\'on usa la variable din\'amica \dsl{*sheet-env*}: durante la
compilaci\'on de cada t\'ermino, la variable de iteraci\'on (\dsl{sheet-var})
se liga temporalmente al nombre de la hoja actual.

\begin{lstlisting}[language=CommonLisp,
  caption={M\'etodo compile-excel-formula para collect-over.}]
(defmethod compile-excel-formula
    ((e clase-xl-expr-collect-over) col-map data-names
     row-num first-row last-row)
  (let* ((groups (groups e))
         (sh-var (sheet-var e))
         (body   (body e))
         (terms  (mapcar
                   (lambda (g)
                     ;; Liga sh-var -> g solo durante este termino
                     (let ((*sheet-env* (acons sh-var g *sheet-env*)))
                       (compile-excel-formula body col-map data-names
                                              row-num first-row last-row)))
                   groups))
         (inner  (format nil "~{~a~^&~}" terms)))
    (format nil "SUBSTITUTE(TRIM(~a),\" \",\",\")" inner)))
\end{lstlisting}

La f\'ormula resultante concatena los $N$ t\'erminos con \texttt{\&} y los
envuelve en \texttt{SUBSTITUTE(TRIM(\ldots))} para eliminar espacios extras
y reemplazarlos con comas, produciendo una lista de grupos separada por comas.

Para el horario de la facultad con 17 grupos, el m\'etodo genera una f\'ormula
que agrega la ocupaci\'on del Aula~1 en el turno~3 del lunes a trav\'es de
las 17 hojas del workbook, sin que el programador haya escrito ning\'un
nombre de hoja ni letra de columna.

% ---------------------------------------------------------------------------
\section{Generaci\'on del programa Python}
\label{sec:generacion-python}
% ---------------------------------------------------------------------------

El generador principal recorre el \'arbol del workbook y emite un programa
Python que usa \texttt{openpyxl} para construir el archivo \xlsx{}.  El
programa generado:

\begin{enumerate}
  \item Crea un workbook openpyxl.
  \item Para cada hoja, crea una hoja de trabajo y la configura.
  \item Para cada tabla de cada regi\'on, escribe las cabeceras y los datos.
  \item Para cada columna calculada, escribe la f\'ormula en cada fila de
    datos (compilada por \dsl{compile-excel-formula} en tiempo de generaci\'on
    del Python).
  \item Para cada regla de estilo condicional, crea una
    \texttt{FormulaRule} con la f\'ormula SUMPRODUCT generada y la a\~nade
    al rango correspondiente.
  \item Para cada expresi\'on fija, resuelve la posici\'on \dsl{nav} y escribe
    la expresi\'on compilada en la celda correspondiente.
  \item Guarda el workbook en el archivo \xlsx{} configurado.
\end{enumerate}

El listado siguiente muestra el fragmento del programa Python generado para
la columna calculada \dsl{defensas} de \dsl{profesores-table}:

\begin{lstlisting}[style=pythonstyle,
  caption={Fragmento del programa Python generado para la columna computed.}]
# Columna: defensas (profesores-table)
ws["M4"] = "=COUNTIF($B$4:$F$6,K4)"
ws["M5"] = "=COUNTIF($B$4:$F$6,K5)"
ws["M6"] = "=COUNTIF($B$4:$F$6,K6)"
ws["M7"] = "=COUNTIF($B$4:$F$6,K7)"
ws["M8"] = "=COUNTIF($B$4:$F$6,K8)"
ws["M9"] = "=COUNTIF($B$4:$F$6,K9)"

# FormulaRule para profesores-table, columna nombre
from openpyxl.formatting.rule import FormulaRule
from openpyxl.styles import PatternFill
red_fill = PatternFill(start_color="FF6666",
                       end_color="FF6666", fill_type="solid")
formula = ("SUMPRODUCT((($B$4:$B$6=$K4)+($C$4:$C$6=$K4)+"
           "($D$4:$D$6=$K4)+($E$4:$E$6=$K4)+($F$4:$F$6=$K4))*"
           "(COUNTIFS($G$4:$G$6,$G$4:$G$6,$H$4:$H$6,$H$4:$H$6)>1))>0")
ws.conditional_formatting.add(
    "K4:K9",
    FormulaRule(formula=[formula], fill=red_fill))
\end{lstlisting}

Las f\'ormulas de las celdas est\'an calculadas en tiempo de generaci\'on del
Python: el programa Python que se ejecuta s\'olo asigna cadenas de texto a las
celdas.  Las referencias relativas dentro de las f\'ormulas las eval\'ua Excel
al abrir el archivo.  La \texttt{FormulaRule}, en cambio, se reevalua
en Excel din\'amicamente cada vez que cambian los datos.
